A new model can be run on only $k$ of $N$ benchmarks: which $k$, how should the other scores be inferred, and how large must $k$ be?
We treat this cold-start problem as Bayesian experimental design on a latent-factor model of the leaderboard score matrix, evaluated under a leakage-safe protocol for sparse, non-stationary leaderboards, with a floor on the error of every selector, a finite-history correction, and two stopping rules.
We test the theory on synthetic matrices with known truth and the selection on MTEB and on generative-LLM leaderboards.
On the complete MTEB block of 103 models and 202 tasks, three benchmarks chosen by the model reduce the error from 1.03 to 0.60 standard deviations, 28\% in MSE below three random ones; but an error of 0.5 takes 20 chosen or 60 random benchmarks, and recovering the ten best history models reaches an F1 of only 0.74 at $k=60$.
The closed-form forecast is accurate with known parameters and optimistic by 10--20\% when fitted on real data; the factor posterior is 22\% less accurate than ExtraTrees at the same selected set, yet the gain from selection survives the change of imputer.
Under temporal shift the nominal 90\% intervals cover 76--80\% of scores on MTEB.
On five generative-LLM leaderboards (6 to 57 benchmarks) the A-optimal criterion beats random choice at the smallest budgets and not at the largest, where the complement-error greedy is never worse; with whole organisations held out the gain at $k=3$ holds on MTEB-A1 and the sparse matrix but not on MTEB-A2, where one organisation carries most of the error, and the advantage at large $k$ does not hold on the sparse matrix.
